\hypertarget{the-standard-template-library-stl-core}{%
\chapter{THE STANDARD TEMPLATE LIBRARY (STL)
CORE}\label{the-standard-template-library-stl-core}}

\hypertarget{c9803-standard-library}{%
\chapter{C++98/03 STANDARD LIBRARY}\label{c9803-standard-library}}

\hypertarget{standard-template-library-stl}{%
\section{Standard Template Library
(STL)}\label{standard-template-library-stl}}

\begin{center}\rule{0.5\linewidth}{0.5pt}\end{center}

\hypertarget{introduction-to-stl}{%
\section{Introduction to STL}\label{introduction-to-stl}}

The Standard Template Library (STL) is a collection of template classes
and functions that provide: - \textbf{Containers} - Data structures to
hold objects (e.g., vectors, lists, maps). - \textbf{Iterators} -
Objects to traverse containers (generalization of pointers). -
\textbf{Algorithms} - Functions to manipulate data (e.g., sorting,
searching). - \textbf{Function Objects (Functors)} - Objects that act
like functions.

\hypertarget{key-advantages}{%
\subsection{Key Advantages}\label{key-advantages}}

\begin{itemize}
\tightlist
\item
  \textbf{Generic Programming}: Write code that works with any data
  type.
\item
  \textbf{Performance}: Heavily optimized implementations.
\item
  \textbf{Reusability}: Standardized components prevent reinventing the
  wheel.
\item
  \textbf{Type Safety}: Templates ensure type correctness at compile
  time.
\end{itemize}

\begin{center}\rule{0.5\linewidth}{0.5pt}\end{center}

\hypertarget{stl-components-overview}{%
\section{STL Components Overview}\label{stl-components-overview}}

\begin{lstlisting}
          STL (Standard Template Library)  
                       |
   -----------------------------------------
   |              |            |           |
CONTAINERS     ITERATORS    ALGORITHMS   FUNCTORS
   |              |            |           |
Sequence       Input        Searching    Predicates
Associative    Output       Sorting      Comparators
Adapters       Forward      Modifying      
               Bidir        Numeric      
               Random       
\end{lstlisting}

\begin{center}\rule{0.5\linewidth}{0.5pt}\end{center}

\hypertarget{containers---complete-reference}{%
\section{CONTAINERS - COMPLETE
REFERENCE}\label{containers---complete-reference}}

\hypertarget{container-characteristics-c98}{%
\subsection{Container Characteristics
(C++98)}\label{container-characteristics-c98}}

\begin{longtable}[]{@{}lllllll@{}}
\toprule
Container & Type & Insert & Delete & Search & Random Access & Memory \\
\midrule
\endhead
\passthrough{\lstinline!vector!} & Sequence & O(n) & O(n) & O(n) & O(1)
& Contiguous \\
\passthrough{\lstinline!list!} & Sequence & O(1) & O(1) & O(n) & - &
Scattered \\
\passthrough{\lstinline!deque!} & Sequence & O(n) & O(n) & O(n) & O(1) &
Chunks \\
\passthrough{\lstinline!map!} & Associative & O(log n) & O(log n) &
O(log n) & - & Tree \\
\passthrough{\lstinline!set!} & Associative & O(log n) & O(log n) &
O(log n) & - & Tree \\
\passthrough{\lstinline!multimap!} & Associative & O(log n) & O(log n) &
O(log n) & - & Tree \\
\passthrough{\lstinline!multiset!} & Associative & O(log n) & O(log n) &
O(log n) & - & Tree \\
\passthrough{\lstinline!priority\_queue!} & Adapter & O(log n) & O(log
n) & - & - & Heap \\
\passthrough{\lstinline!queue!} & Adapter & O(1) & O(1) & - & - & - \\
\passthrough{\lstinline!stack!} & Adapter & O(1) & O(1) & - & - & - \\
\bottomrule
\end{longtable}

\begin{center}\rule{0.5\linewidth}{0.5pt}\end{center}

\hypertarget{vector---dynamic-array}{%
\section{1.1 VECTOR - Dynamic Array}\label{vector---dynamic-array}}

\hypertarget{what-is-vector}{%
\subsection{What is Vector?}\label{what-is-vector}}

A dynamic array that grows automatically. Use this as your default
container.

\hypertarget{declaration-initialization}{%
\subsection{Declaration \&
Initialization}\label{declaration-initialization}}

\begin{lstlisting}[language={C++}]
#include <vector>
using namespace std;

// Empty vector
vector<int> v1;

// Vector with initial size (10 elements initialized to 0)
vector<int> v2(10);

// Vector with initial size and value (5 elements, all 10)
vector<int> v3(5, 10);

// Copy constructor
vector<int> v4(v3);

// From array (using pointers)
int arr[] = {1, 2, 3, 4, 5};
vector<int> v5(arr, arr + 5); 
\end{lstlisting}

\hypertarget{accessing-elements}{%
\subsection{Accessing Elements}\label{accessing-elements}}

\begin{lstlisting}[language={C++}]
vector<int> v;
v.push_back(10);
v.push_back(20);

// 1. Using operator[] (No bounds check)
cout << v[0] << "\n";  // 10

// 2. Using at() (With bounds check, throws std::out_of_range)
cout << v.at(1) << "\n"; // 20

// 3. Front and back
cout << v.front() << "\n"; // 10
cout << v.back() << "\n";  // 20
\end{lstlisting}

\hypertarget{modifying-elements}{%
\subsection{Modifying Elements}\label{modifying-elements}}

\begin{lstlisting}[language={C++}]
vector<int> v;

// Adding elements
v.push_back(10);
v.push_back(20);
v.push_back(30); // {10, 20, 30}

// Inserting elements (Iterators required)
v.insert(v.begin() + 1, 15); // {10, 15, 20, 30}

// Removing elements
v.pop_back();           // Removes 30
v.erase(v.begin() + 1); // Removes 15

// Clearing
v.clear();              // Empty
\end{lstlisting}

\hypertarget{size-capacity}{%
\subsection{Size \& Capacity}\label{size-capacity}}

\begin{lstlisting}[language={C++}]
vector<int> v(10);

cout << v.size() << "\n";      // 10
cout << v.capacity() << "\n";  // >= 10
cout << v.empty() << "\n";     // 0 (false)

// Reserve space (Optimization to avoid reallocations)
v.reserve(100); 

// Resize (Changes size, fills new elements with default or value)
v.resize(20);    // Size becomes 20
v.resize(5);     // Size becomes 5, extra elements destroyed
\end{lstlisting}

\hypertarget{iterating-through-vector}{%
\subsection{Iterating Through Vector}\label{iterating-through-vector}}

\begin{lstlisting}[language={C++}]
vector<int> v;
v.push_back(10); v.push_back(20); v.push_back(30);

// 1. Index loop
for (size_t i = 0; i < v.size(); ++i) {
    cout << v[i] << " ";
}

// 2. Iterator loop (Recommended)
for (vector<int>::iterator it = v.begin(); it != v.end(); ++it) {
    cout << *it << " ";
}

// 3. Reverse Iterator
for (vector<int>::reverse_iterator rit = v.rbegin(); rit != v.rend(); ++rit) {
    cout << *rit << " ";
}
\end{lstlisting}

\begin{center}\rule{0.5\linewidth}{0.5pt}\end{center}

\hypertarget{deque---double-ended-queue}{%
\section{1.2 DEQUE - Double Ended
Queue}\label{deque---double-ended-queue}}

\hypertarget{what-is-deque}{%
\subsection{What is Deque?}\label{what-is-deque}}

Fast insertion/deletion at both ends. Unlike vector, storage is not
guaranteed to be contiguous.

\begin{lstlisting}[language={C++}]
#include <deque>
using namespace std;

deque<int> dq;

dq.push_back(10);
dq.push_front(5);  // {5, 10}

dq.pop_back();     // {5}
dq.pop_front();    // {}

// Supports random access
dq.push_back(100);
cout << dq[0] << "\n";
\end{lstlisting}

\begin{center}\rule{0.5\linewidth}{0.5pt}\end{center}

\hypertarget{list---doubly-linked-list}{%
\section{1.3 LIST - Doubly Linked
List}\label{list---doubly-linked-list}}

\hypertarget{what-is-list}{%
\subsection{What is List?}\label{what-is-list}}

Efficient insertion/deletion anywhere (O(1) if iterator is known). No
random access.

\begin{lstlisting}[language={C++}]
#include <list>
using namespace std;

list<int> lst;
lst.push_back(10);
lst.push_front(5);

// Insert at position
list<int>::iterator it = lst.begin();
++it; // Move to second position
lst.insert(it, 7); // {5, 7, 10}

// Remove
lst.remove(7); // Remove all elements with value 7

// Unique (Removes consecutive duplicates)
lst.push_back(10);
lst.unique(); // {5, 10} (second 10 removed)

// Sort
lst.sort(); // Internal sort (std::sort doesn't work on lists)
\end{lstlisting}

\begin{center}\rule{0.5\linewidth}{0.5pt}\end{center}

\hypertarget{map---sorted-key-value-pairs}{%
\section{1.4 MAP - Sorted Key-Value
Pairs}\label{map---sorted-key-value-pairs}}

\hypertarget{what-is-map}{%
\subsection{What is Map?}\label{what-is-map}}

Associative container storing key-value pairs sorted by key. Implemented
as a Balanced BST (usually Red-Black Tree).

\begin{lstlisting}[language={C++}]
#include <map>
#include <string>
using namespace std;

map<string, int> ages;

// Insertion
ages["Alice"] = 30;
ages["Bob"] = 25;
ages.insert(make_pair("Charlie", 28));

// Access / Search
cout << ages["Alice"] << "\n"; // 30

map<string, int>::iterator it = ages.find("Bob");
if (it != ages.end()) {
    cout << "Bob is " << it->second << " years old.\n";
}

// Iteration (Sorted by key)
for (it = ages.begin(); it != ages.end(); ++it) {
    cout << it->first << ": " << it->second << "\n";
}
\end{lstlisting}

\begin{center}\rule{0.5\linewidth}{0.5pt}\end{center}

\hypertarget{set---sorted-unique-keys}{%
\section{1.5 SET - Sorted Unique Keys}\label{set---sorted-unique-keys}}

\hypertarget{what-is-set}{%
\subsection{What is Set?}\label{what-is-set}}

Stores unique elements sorted by value.

\begin{lstlisting}[language={C++}]
#include <set>
using namespace std;

set<int> s;
s.insert(10);
s.insert(5);
s.insert(10); // Duplicate ignored

// {5, 10}

if (s.find(5) != s.end()) {
    cout << "5 is present.\n";
}

// Range erase
s.erase(s.begin()); // Removes 5
\end{lstlisting}

\begin{center}\rule{0.5\linewidth}{0.5pt}\end{center}

\hypertarget{multimap-multiset}{%
\section{1.6 MULTIMAP \& MULTISET}\label{multimap-multiset}}

Allow duplicate keys (multimap) or duplicate values (multiset).

\begin{lstlisting}[language={C++}]
#include <map>
using namespace std;

multimap<string, int> scores;
scores.insert(make_pair("Alice", 90));
scores.insert(make_pair("Alice", 85)); // Allowed

// Finding all values for a key
pair<multimap<string, int>::iterator, multimap<string, int>::iterator> range;
range = scores.equal_range("Alice");

for (multimap<string, int>::iterator it = range.first; it != range.second; ++it) {
    cout << it->second << " ";
}
// Output: 90 85
\end{lstlisting}

\begin{center}\rule{0.5\linewidth}{0.5pt}\end{center}

\hypertarget{stack-adapter}{%
\section{1.7 STACK (Adapter)}\label{stack-adapter}}

LIFO (Last In, First Out). Wraps \passthrough{\lstinline!deque!}
(default), \passthrough{\lstinline!vector!}, or
\passthrough{\lstinline!list!}.

\begin{lstlisting}[language={C++}]
#include <stack>
using namespace std;

stack<int> st;
st.push(10);
st.push(20);

cout << st.top() << "\n"; // 20
st.pop();                 // Removes 20
\end{lstlisting}

\begin{center}\rule{0.5\linewidth}{0.5pt}\end{center}

\hypertarget{queue-adapter}{%
\section{1.8 QUEUE (Adapter)}\label{queue-adapter}}

FIFO (First In, First Out). Wraps \passthrough{\lstinline!deque!}
(default) or \passthrough{\lstinline!list!}.

\begin{lstlisting}[language={C++}]
#include <queue>
using namespace std;

queue<int> q;
q.push(10);
q.push(20);

cout << q.front() << "\n"; // 10
q.pop();                   // Removes 10
\end{lstlisting}

\begin{center}\rule{0.5\linewidth}{0.5pt}\end{center}

\hypertarget{priority_queue-adapter}{%
\section{1.9 PRIORITY\_QUEUE (Adapter)}\label{priority_queue-adapter}}

Sorted queue (Heap). Max element is always at
\passthrough{\lstinline!top()!}.

\begin{lstlisting}[language={C++}]
#include <queue>
#include <vector>
#include <functional> // for greater<int>
using namespace std;

// Max Heap (Default)
priority_queue<int> pq;
pq.push(10);
pq.push(30);
pq.push(20);

cout << pq.top() << "\n"; // 30

// Min Heap
priority_queue<int, vector<int>, greater<int> > min_pq;
min_pq.push(10);
min_pq.push(30);

cout << min_pq.top() << "\n"; // 10
\end{lstlisting}

\begin{center}\rule{0.5\linewidth}{0.5pt}\end{center}

\hypertarget{iterators}{%
\section{ITERATORS}\label{iterators}}

Iterators are the glue between Containers and Algorithms.

\hypertarget{categories}{%
\subsection{Categories}\label{categories}}

\begin{enumerate}
\def\labelenumi{\arabic{enumi}.}
\tightlist
\item
  \textbf{Input}: Read-only, single pass (e.g.,
  \passthrough{\lstinline!istream\_iterator!}).
\item
  \textbf{Output}: Write-only, single pass (e.g.,
  \passthrough{\lstinline!ostream\_iterator!}).
\item
  \textbf{Forward}: Read/Write, multi-pass (e.g.,
  \passthrough{\lstinline!slist!} - non-std).
\item
  \textbf{Bidirectional}: Forward + Backward (e.g.,
  \passthrough{\lstinline!list!}, \passthrough{\lstinline!map!},
  \passthrough{\lstinline!set!}).
\item
  \textbf{Random Access}: O(1) jump (e.g.,
  \passthrough{\lstinline!vector!}, \passthrough{\lstinline!deque!},
  arrays).
\end{enumerate}

\hypertarget{operations}{%
\subsection{Operations}\label{operations}}

\begin{itemize}
\tightlist
\item
  \passthrough{\lstinline!*it!}: Dereference.
\item
  \passthrough{\lstinline!++it!}: Advance.
\item
  \passthrough{\lstinline!--it!}: Retreat (Bidirectional+).
\item
  \passthrough{\lstinline!it + n!}: Jump (Random Access only).
\item
  \passthrough{\lstinline!it1 - it2!}: Distance (Random Access only).
\end{itemize}

\begin{lstlisting}[language={C++}]
vector<int> v(5, 1); 
vector<int>::iterator it = v.begin();
advance(it, 2); // Move 2 steps
cout << distance(v.begin(), it); // 2
\end{lstlisting}

\begin{center}\rule{0.5\linewidth}{0.5pt}\end{center}

\hypertarget{algorithms}{%
\section{ALGORITHMS}\label{algorithms}}

Found in \passthrough{\lstinline!<algorithm>!} and
\passthrough{\lstinline!<numeric>!}. They work on iterator ranges
\passthrough{\lstinline![first, last)!}.

\hypertarget{non-modifying}{%
\subsection{1. Non-Modifying}\label{non-modifying}}

\begin{itemize}
\tightlist
\item
  \passthrough{\lstinline!find(begin, end, val)!}
\item
  \passthrough{\lstinline!count(begin, end, val)!}
\item
  \passthrough{\lstinline!equal(b1, e1, b2)!}
\item
  \passthrough{\lstinline!search(b1, e1, b2, e2)!}
\end{itemize}

\hypertarget{modifying}{%
\subsection{2. Modifying}\label{modifying}}

\begin{itemize}
\tightlist
\item
  \passthrough{\lstinline!copy(b1, e1, b2)!}
\item
  \passthrough{\lstinline!transform(b1, e1, out, op)!}
\item
  \passthrough{\lstinline!replace(b1, e1, old, new)!}
\item
  \passthrough{\lstinline!fill(b, e, val)!}
\item
  \passthrough{\lstinline!swap(a, b)!}
\item
  \passthrough{\lstinline!reverse(b, e)!}
\item
  \passthrough{\lstinline!rotate(b, mid, e)!}
\item
  \passthrough{\lstinline!random\_shuffle(b, e)!}
\end{itemize}

\hypertarget{sorting}{%
\subsection{3. Sorting}\label{sorting}}

\begin{itemize}
\tightlist
\item
  \passthrough{\lstinline!sort(b, e)!}: O(N log N).
\item
  \passthrough{\lstinline!stable\_sort(b, e)!}: Preserves order of equal
  elements.
\item
  \passthrough{\lstinline!partial\_sort(b, mid, e)!}: Top K elements.
\end{itemize}

\hypertarget{binary-search-on-sorted-ranges}{%
\subsection{4. Binary Search (On Sorted
Ranges)}\label{binary-search-on-sorted-ranges}}

\begin{itemize}
\tightlist
\item
  \passthrough{\lstinline!binary\_search(b, e, val)!}: Returns bool.
\item
  \passthrough{\lstinline!lower\_bound(b, e, val)!}: First element
  \textgreater= val.
\item
  \passthrough{\lstinline!upper\_bound(b, e, val)!}: First element
  \textgreater{} val.
\end{itemize}

\hypertarget{set-operations-on-sorted-ranges}{%
\subsection{5. Set Operations (On Sorted
Ranges)}\label{set-operations-on-sorted-ranges}}

\begin{itemize}
\tightlist
\item
  \passthrough{\lstinline!set\_union!},
  \passthrough{\lstinline!set\_intersection!},
  \passthrough{\lstinline!set\_difference!}.
\end{itemize}

\hypertarget{numeric}{%
\subsection{6. Numeric ()}\label{numeric}}

\begin{itemize}
\tightlist
\item
  \passthrough{\lstinline!accumulate(b, e, init)!}: Sum.
\item
  \passthrough{\lstinline!inner\_product!}: Dot product.
\item
  \passthrough{\lstinline!adjacent\_difference!}.
\end{itemize}

\hypertarget{example-sort-and-find}{%
\subsection{Example: Sort and Find}\label{example-sort-and-find}}

\begin{lstlisting}[language={C++}]
#include <algorithm>
#include <vector>
#include <iostream>
using namespace std;

bool descending(int a, int b) {
    return a > b;
}

int main() {
    int arr[] = {5, 2, 9, 1, 5, 6};
    vector<int> v(arr, arr + 6);

    // Sort ascending
    sort(v.begin(), v.end());

    // Binary Search
    if (binary_search(v.begin(), v.end(), 9)) {
        cout << "Found 9\n";
    }

    // Sort descending with predicate
    sort(v.begin(), v.end(), descending);

    return 0;
}
\end{lstlisting}

\begin{center}\rule{0.5\linewidth}{0.5pt}\end{center}

\hypertarget{strings-stdstring}{%
\section{STRINGS (std::string)}\label{strings-stdstring}}

A specialization of \passthrough{\lstinline!basic\_string<char>!}.
Replaces C-style \passthrough{\lstinline!char*!} strings.

\begin{lstlisting}[language={C++}]
#include <string>
#include <iostream>
using namespace std;

string s = "Hello";
s += " World"; // Concatenation

// Substring
string sub = s.substr(0, 5); // "Hello"

// Find
size_t pos = s.find("World");
if (pos != string::npos) {
    cout << "Found at " << pos << "\n";
}

// C-string compatibility
const char* cstr = s.c_str();
\end{lstlisting}

\begin{center}\rule{0.5\linewidth}{0.5pt}\end{center}

\hypertarget{streams-iostream}{%
\section{STREAMS (IOSTREAM)}\label{streams-iostream}}

\hypertarget{file-io}{%
\subsection{File I/O ()}\label{file-io}}

\begin{lstlisting}[language={C++}]
#include <fstream>
#include <iostream>
using namespace std;

int main() {
    // Write
    ofstream out("test.txt");
    if (out) {
        out << "Line 1" << endl;
        out << 123 << endl;
    }
    out.close();

    // Read
    ifstream in("test.txt");
    string line;
    int num;
    
    if (in >> line >> num) { // Reads "Line" then fails on "1" vs int? No. 
                             // "Line" goes to line. "1" goes to num? 
                             // Need careful parsing.
    }
    // Better: getline
    getline(in, line);
    
    return 0;
}
\end{lstlisting}

\hypertarget{string-streams}{%
\subsection{String Streams ()}\label{string-streams}}

Useful for parsing strings or formatting.

\begin{lstlisting}[language={C++}]
#include <sstream>
using namespace std;

// Int to String
int x = 42;
stringstream ss;
ss << x;
string s = ss.str();

// String to Int
string s2 = "100";
stringstream ss2(s2);
int y;
ss2 >> y;
\end{lstlisting}

\begin{center}\rule{0.5\linewidth}{0.5pt}\end{center}

\hypertarget{functors-function-objects}{%
\section{FUNCTORS (Function Objects)}\label{functors-function-objects}}

Classes that overload \passthrough{\lstinline!operator()!}. Used by
algorithms.

\begin{lstlisting}[language={C++}]
struct AddX {
    int x;
    AddX(int val) : x(val) {}
    
    int operator()(int y) const {
        return x + y;
    }
};

// Usage with transform
vector<int> v(5, 1); // {1, 1, 1, 1, 1}
transform(v.begin(), v.end(), v.begin(), AddX(10));
// v is now {11, 11, 11, 11, 11}
\end{lstlisting}

This covers the foundational C++98 Standard Library components necessary
for mastery.

\hypertarget{advanced-strings}{%
\chapter{ADVANCED STRINGS}\label{advanced-strings}}

\hypertarget{string-manipulation}{%
\section{5.1 String Manipulation}\label{string-manipulation}}

\begin{lstlisting}[language={C++}]
#include <iostream>
#include <string>
#include <cstring>
using namespace std;

int main() {
    string s = "Hello World";
    
    // Length and capacity
    cout << "Length: " << s.length() << endl;
    cout << "Capacity: " << s.capacity() << endl;
    
    // Access characters
    cout << "First char: " << s[0] << endl;
    cout << "Last char: " << s[s.length() - 1] << endl;
    
    // Finding substrings
    size_t pos = s.find("World");
    if (pos != string::npos) {
        cout << "Found at position: " << pos << endl;
    }
    
    // Replace
    s.replace(6, 5, "C++");
    cout << s << endl;  // Hello C++
    
    // Insert
    s.insert(5, " there");
    cout << s << endl;  // Hello there C++
    
    // Erase
    s.erase(5, 6);
    cout << s << endl;  // Hello C++
    
    // Substring
    cout << s.substr(0, 5) << endl;  // Hello
    
    // Reverse
    reverse(s.begin(), s.end());
    cout << s << endl;  // ++C olleH
    
    return 0;
}
\end{lstlisting}

\hypertarget{string-conversion}{%
\section{5.2 String Conversion}\label{string-conversion}}

\begin{lstlisting}[language={C++}]
#include <iostream>
#include <string>
#include <sstream>
#include <cstdlib>
using namespace std;

int main() {
    // String to number (C style)
    string s1 = "42";
    int num = atoi(s1.c_str());
    cout << num << endl;
    
    string s2 = "3.14";
    double dbl = atof(s2.c_str());
    cout << dbl << endl;
    
    // Number to string (using stringstream)
    stringstream ss;
    ss << 42 << " " << 3.14 << " " << true;
    string result = ss.str();
    cout << result << endl;  // 42 3.14 1
    
    // Reverse conversion
    stringstream ss2("100 200 300");
    int a, b, c;
    ss2 >> a >> b >> c;
    cout << a << " " << b << " " << c << endl;  // 100 200 300
    
    return 0;
}
\end{lstlisting}

\hypertarget{string-tokenization}{%
\section{5.3 String Tokenization}\label{string-tokenization}}

\begin{lstlisting}[language={C++}]
#include <iostream>
#include <string>
#include <cstring>
using namespace std;

int main() {
    string line = "apple,banana,orange,grape";
    
    // Using stringstream and getline
    stringstream ss(line);
    string token;
    
    while (getline(ss, token, ',')) {
        cout << token << endl;
    }
    // Output: apple, banana, orange, grape
    
    // Using strtok (C style)
    char str[] = "hello world how are you";
    char* ptr = strtok(str, " ");
    
    while (ptr != NULL) {
        cout << ptr << endl;
        ptr = strtok(NULL, " ");
    }
    // Output: hello, world, how, are, you
    
    return 0;
}
\end{lstlisting}

\begin{center}\rule{0.5\linewidth}{0.5pt}\end{center}

\hypertarget{file-io-advanced}{%
\chapter{FILE I/O ADVANCED}\label{file-io-advanced}}

\hypertarget{binary-file-io}{%
\section{14.1 Binary File I/O}\label{binary-file-io}}

\begin{lstlisting}[language={C++}]
#include <iostream>
#include <fstream>
using namespace std;

int main() {
    // Write binary data
    ofstream outfile("data.bin", ios::binary);
    
    int numbers[] = {10, 20, 30, 40, 50};
    outfile.write((char*)numbers, sizeof(numbers));
    outfile.close();
    
    // Read binary data
    ifstream infile("data.bin", ios::binary);
    
    int buffer[5];
    infile.read((char*)buffer, sizeof(buffer));
    
    for (int i = 0; i < 5; i++) {
        cout << buffer[i] << " ";
    }
    cout << endl;
    
    infile.close();
    
    return 0;
}
\end{lstlisting}

\hypertarget{stream-positioning}{%
\section{14.2 Stream Positioning}\label{stream-positioning}}

\begin{lstlisting}[language={C++}]
#include <iostream>
#include <fstream>
using namespace std;

int main() {
    // Write to file
    ofstream outfile("test.txt");
    outfile << "0123456789";
    outfile.close();
    
    // Read with positioning
    ifstream infile("test.txt");
    
    // Tell position
    cout << "Current position: " << infile.tellg() << endl;
    
    // Seek to position
    infile.seekg(5);
    char c;
    infile.get(c);
    cout << "Character at position 5: " << c << endl;  // '5'
    
    // Seek from end
    infile.seekg(-3, ios::end);
    infile.get(c);
    cout << "Third from end: " << c << endl;  // '7'
    
    infile.close();
    
    return 0;
}
\end{lstlisting}

\begin{center}\rule{0.5\linewidth}{0.5pt}\end{center}
